\section{Introduction}
\label{sec:introduction}

Electroencephalography (EEG) is a physiological signal that records brain neural activity non-invasively, and it is widely used in cognitive neuroscience, clinical neurology, and brain--computer interfaces (BCIs)~\cite{chen2022openworld}. Owing to its high temporal resolution and convenient acquisition, EEG plays a key role in neural-function monitoring, computer-aided diagnosis, and assessment of states of consciousness. 
Nevertheless, EEG has an intrinsically low signal-to-noise ratio and is easily contaminated by a variety of artifacts and environmental noise, including ocular (EOG) and electromyographic (EMG) activity, power-line interference (e.g., the 50/60\,Hz line frequency), and poor electrode contact. Such contamination severely distorts the underlying neural information and degrades the accuracy of downstream analyses~\cite{zhang2021eegdenoisenet,yin2022frequency}.

To suppress noise and extract reliable neural features, EEG denoising has long been a core problem in the field. Traditional approaches---most prominently band-pass filtering and Independent Component Analysis (ICA)---perform well in specific settings but suffer from two key bottlenecks. First, they are highly sensitive to parameter choices and to assumptions about the noise distribution, which makes them difficult to adapt to the dynamic and complex conditions of real-world acquisition~\cite{ghosh2019eyeblink}. Second, they
largely ignore inter-individual differences in cranial anatomy, cognitive state, and acquisition conditions, so their generalization deteriorates in cross-subject scenarios~\cite{jiao2018sparse}.

In recent years, deep learning has provided new tools for modeling and denoising EEG signals. Following the success of generative adversarial networks (GANs) and variational autoencoders (VAEs) in image denoising and audio reconstruction, considerable effort has been devoted to applying such models to neural-signal processing~\cite{yin2023gan}. However, these models typically require large amounts of identically distributed training data; faced with the high heterogeneity and non-stationarity of EEG, they still struggle with generalization and training stability.

Denoising Diffusion Probabilistic Models (DDPMs) are a recently emerged generative framework that has achieved major breakthroughs in image and audio generation~\cite{ho2020ddpm}, thanks to progressive reconstruction, stable training, and a strong ability to model complex distributions. The core idea is to gradually perturb data into noise and then recover clean samples through a reverse diffusion process. This layer-by-layer formulation is naturally aligned with the need to progressively purify EEG signals and restore their structure. Although diffusion models have been explored for the reconstruction of other physiological signals such as the electrocardiogram~\cite{zhao2023ecg,goldberger2000physionet}, DDPM-based modeling that accounts for cross-subject individual differences in EEG remains at an early stage and has not been systematically studied~\cite{zeng2023dmre2i}.

This paper presents a new strategy for cross-subject EEG denoising, the \emph{Subject-Aware Denoising Diffusion Probabilistic Model} (\saddpm{}). To the best of our knowledge, this is among the first attempts to incorporate a subject-aware mechanism into the DDPM denoising pipeline, endowing the model with the ability to recognize and adapt to inter-individual domain differences. During modeling, a distinct embedding vector is learned for each subject and is injected as a conditioning signal at multiple layers of the reverse process, so that the model can dynamically adjust its noise-estimation trajectory according to the characteristics of each subject, so as to pursue more robust and transferable denoising. The central idea is to combine the expressive power of generative modeling with the generalization afforded by an individual-aware mechanism, yielding a denoising model that is physiologically motivated and has practical potential for handling the structural artifacts and individual variability inherent in EEG. In short, this work builds a DDPM-based, domain-specific EEG denoising method that explicitly targets the weak cross-subject generalization of existing approaches. From the perspective of autonomous and adaptive systems, an EEG pipeline that must remain effective for new users is itself a learning-based system facing distribution shift, and our subject-aware diffusion model instantiates the emerging use of generative AI as a mechanism for such adaptation~\cite{li2024genai,casimiro2024selfadapting}. Our main contributions are summarized as follows.

\begin{itemize}
  \item We design and implement a domain-specific denoising diffusion probabilistic model (\saddpm{}) that exploits the step-by-step denoising of diffusion models to effectively distinguish the noise and signal characteristics of different subjects.

  \item Through a transductive, label-free target-subject adaptation strategy, SADDPM estimates a deployment subject's embedding from unlabeled target EEG and adapts the reverse denoising trajectory without using target task labels, clean targets, or artifact annotations.

  \item We construct a complete pipeline---covering data processing, model implementation, and evaluation---and conduct systematic experiments on the BCI Competition~IV-2a dataset, validating the effectiveness of the proposed method under complex noise conditions and in cross-subject tasks.

  \item We complement the downstream-classification study with a paired ground-truth denoising evaluation under the EEGdenoiseNet protocol using a conditional variant, \saddpmcond{}. The results show competitive single-channel ocular-artifact removal, high-fidelity multi-channel denoising, and functionally relevant subject conditioning when artifacts are subject-specific.
\end{itemize}

The remainder of this paper is organized as follows.
Section~\ref{sec:related} reviews related work on EEG denoising, subject variability, and diffusion-based signal modeling.
Section~\ref{sec:method} details the proposed \saddpm{}, including its diffusion formulation, subject-aware conditioning, network architecture, and training objective. Section~\ref{sec:experiments} describes the dataset, preprocessing, experimental protocol, and downstream results; Section~\ref{sec:exp_phase2} reports the paired ground-truth denoising study. Section~\ref{sec:conclusion} concludes the paper and discusses future directions.
